% =====================================================
% 题型：抛物线焦点距离选择题 (q_parabola_focus_distance.tex)
% 知识板块：抛物线
% 主思维：数形结合与几何直观
% 副思维：定义域/定义优先、化归转化
% 错因标签：抛物线焦点公式记混（开口向右 vs 开口向上）
% 讲义用途：解析几何基础
% =====================================================
% 难度：★ (基础)
% =====================================================
\newcommand{\qParabolaFocusDistanceStem}{已知抛物线 \(C_1: y^2 = 2p_1x \ (p_1 > 0)\) 和 \(C_2: x^2 = 2p_2y \ (p_2 > 0)\) 均经过点 \((4, 8)\)，则 \(C_1\) 的焦点与 \(C_2\) 的焦点之间的距离为 \hfill （\quad）}

\newcommand{\qParabolaFocusDistanceOptions}{%
  \mcoptions[2]{\(12\)}{\(4\sqrt{5}\)}{\(6\)}{\(\dfrac{\sqrt{65}}{2}\)}%
}

\newcommand{\qParabolaFocusDistanceAnswer}{D}

\newcommand{\qParabolaFocusDistanceSolution}{%
  【解题思路】\\
  第一步，将公共点 \((4, 8)\) 的坐标分别代入两条抛物线方程，求出参数 \(p_1\) 和 \(p_2\)；第二步，写出两条不同开口方向抛物线的焦点坐标公式（注意开口向右的焦点在 \(x\) 轴上，开口向上的焦点在 \(y\) 轴上），从而得出对应的焦点坐标 \(F_1\) 和 \(F_2\)；第三步，利用两点间距离公式求出两焦点之间的距离。\\
  【解析计算】\\
  对于第一条抛物线 \(C_1: y^2 = 2p_1x\)，将点 \((4, 8)\) 代入得：\\
  \(8^2 = 2p_1 \cdot 4 \implies 64 = 8p_1 \implies p_1 = 8\)。\\
  对于抛物线 \(y^2 = 2px\)，其焦点坐标为 \(\left(\dfrac{p}{2}, 0\right)\)，因此 \(C_1\) 的焦点为 \(F_1(4, 0)\)。\\
  对于第二条抛物线 \(C_2: x^2 = 2p_2y\)，将点 \((4, 8)\) 代入得：\\
  \(4^2 = 2p_2 \cdot 8 \implies 16 = 16p_2 \implies p_2 = 1\)。\\
  对于抛物线 \(x^2 = 2py\)，其焦点坐标为 \(\left(0, \dfrac{p}{2}\right)\)，因此 \(C_2\) 的焦点为 \(F_2\left(0, \dfrac{1}{2}\right)\)。\\
  根据两点间距离公式，两焦点之间的距离为：\\
  \(d = \sqrt{(4 - 0)^2 + \left(0 - \dfrac{1}{2}\right)^2} = \sqrt{16 + \dfrac{1}{4}} = \sqrt{\dfrac{65}{4}} = \dfrac{\sqrt{65}}{2}\)。\\
  故选 D。%
}

\newcommand{\qParabolaFocusDistanceDiagnosis}{%
  学生极易混淆抛物线的焦点公式。由于两条抛物线开口方向不同（一个向右，一个向上），其焦点坐标形式分别为 \(\left(\dfrac{p}{2}, 0\right)\) 和 \(\left(0, \dfrac{p}{2}\right)\)。如果套错公式（例如将 \(x^2=2py\) 的焦点写在 \(x\) 轴上），会导致计算失误。%
}

\newcommand{\qParabolaFocusDistanceTeachingNotes}{%
  讲评时可以画图帮助学生建立几何直观。对比强调四种标准抛物线的焦点与准线方程，要求学生做题时先判断开口方向，再确定坐标位置。%
}
